% Manual del equipo de Micros Tempo, versión alfa.
%
% Se compila con:   cd equipo && latexmk -pdf manual-del-equipo.tex
% Usa el preámbulo del plan (plan-financiero/latex/preambulo.tex): si cambia
% la paleta o un estilo, cambia en los dos documentos.
% Las metas y los sueldos son los del Plan del primer año, versión alfa.
\documentclass[a4paper]{article}
% Preámbulo común del plan y del manual del equipo (plan-financiero/latex).
% La paleta es la del producto, la misma de manual/manual.tex y legal/preambulo.tex.
\usepackage[T1]{fontenc}
\usepackage[utf8]{inputenc}
\usepackage[spanish,es-noshorthands,es-tabla,es-nodecimaldot]{babel}
\usepackage[scaled=0.92]{helvet}
\renewcommand{\familydefault}{\sfdefault}
\usepackage{textcomp}
\usepackage{sansmath}
\AtBeginDocument{\sansmath}
\usepackage{amssymb}
\DeclareUnicodeCharacter{2212}{\textminus}
\DeclareUnicodeCharacter{00BD}{\textonehalf}
\usepackage[a4paper,top=2cm,bottom=2.2cm,left=2.2cm,right=2.2cm,headheight=14pt]{geometry}
\usepackage[table]{xcolor}
\usepackage{microtype}
\usepackage{enumitem}
\usepackage{booktabs}
\usepackage{tabularx}
\usepackage{array}
\usepackage{multirow}
\usepackage{fancyhdr}
\usepackage{titlesec}
\usepackage{lastpage}
\usepackage{needspace}
\usepackage{etoolbox}
\usepackage{pdflscape}
\usepackage{afterpage}
\usepackage{caption}
\usepackage{float}
\usepackage[most]{tcolorbox}
\usepackage{pgfplots}
\pgfplotsset{compat=1.18}
\usepgfplotslibrary{fillbetween}
\usetikzlibrary{patterns}
\usepackage[hidelinks]{hyperref}

\definecolor{tinta}{HTML}{0B2239}
\definecolor{tinta2}{HTML}{2A4864}
\definecolor{tenue}{HTML}{496A86}
\definecolor{marca}{HTML}{1668C4}
\definecolor{verde}{HTML}{1F8A4F}
\definecolor{ambar}{HTML}{A67300}
\definecolor{naranja}{HTML}{D9622B}
\definecolor{rojo}{HTML}{C2001D}
\definecolor{rojosuave}{HTML}{D9827A}
\definecolor{fondo}{HTML}{F4F7FB}
\definecolor{linea}{HTML}{C3D4E4}
\definecolor{rejilla}{HTML}{E3EAF2}
\definecolor{ambarfondo}{HTML}{FDF6E4}
\definecolor{ambarborde}{HTML}{E4C88A}
\definecolor{azulfondo}{HTML}{E7F0FB}
\definecolor{azulborde}{HTML}{C9DDF4}
\definecolor{verdefondo}{HTML}{E9F7EF}

\color{tinta}
\setlength{\parindent}{0pt}
\setlength{\parskip}{5pt plus 2pt}
\linespread{1.08}
\setlist{itemsep=2pt,topsep=3pt,leftmargin=14pt}

% Ningún párrafo se parte entre dos páginas, y ningún título queda al pie
% con su texto en la página siguiente: antes de cada título se pide lugar
% para el título y al menos un párrafo corto.
\interlinepenalty=10000
\clubpenalty=10000
\widowpenalty=10000
\raggedbottom
\titleformat{\section}{\color{tinta}\Large\bfseries}{\thesection.}{0.5em}{}
\titleformat{\subsection}{\color{tinta2}\large\bfseries}{}{0em}{}
\titleformat{\subsubsection}{\color{tinta2}\normalsize\bfseries}{}{0em}{}
\titlespacing*{\section}{0pt}{14pt}{6pt}
\titlespacing*{\subsection}{0pt}{12pt}{3pt}
\titlespacing*{\subsubsection}{0pt}{8pt}{2pt}
\pretocmd{\section}{\needspace{14\baselineskip}}{}{}
\pretocmd{\subsection}{\needspace{10\baselineskip}}{}{}
\pretocmd{\subsubsection}{\needspace{7\baselineskip}}{}{}

\renewcommand{\arraystretch}{1.25}
\arrayrulecolor{linea}
\captionsetup{font={small,color=tenue},labelformat=empty,skip=4pt,justification=raggedright,singlelinecheck=false}

\newtcolorbox{nota}[1][]{colback=azulfondo,colframe=azulborde,boxrule=0.6pt,arc=3pt,
  left=8pt,right=8pt,top=6pt,bottom=6pt,fontupper=\small\color{tinta2},breakable=false,#1}
\newtcolorbox{aviso}[1][]{colback=ambarfondo,colframe=ambarborde,boxrule=0.6pt,arc=3pt,
  left=8pt,right=8pt,top=6pt,bottom=6pt,fontupper=\small,breakable=false,#1}
\newtcolorbox{rol}[1]{colback=white,colframe=linea,boxrule=0.6pt,arc=4pt,
  left=9pt,right=9pt,top=7pt,bottom=5pt,breakable=false,
  title={#1},coltitle=tinta,colbacktitle=fondo,fonttitle=\bfseries,toptitle=3pt,bottomtitle=3pt}

% Cifra destacada: valor grande y su explicación
\newcommand{\cifra}[3][tinta]{%
  \begin{tcolorbox}[colback=white,colframe=linea,boxrule=0.6pt,arc=4pt,left=7pt,right=7pt,top=5pt,bottom=5pt,height=2.45cm,valign=top]
  {\color{#1}\fontsize{15}{18}\selectfont\bfseries #2}\par\vspace{2pt}
  {\footnotesize\color{tenue}#3}
  \end{tcolorbox}}

\newcommand{\documento}[1]{%
  \pagestyle{fancy}\fancyhf{}%
  \renewcommand{\headrulewidth}{0pt}%
  \fancyfoot[L]{\scriptsize\color{tenue}#1}%
  \fancyfoot[R]{\scriptsize\color{tenue}Página \thepage\ de \pageref*{LastPage}}%
}

% Estilo común de los gráficos
\pgfplotsset{
  casa/.style={
    width=\linewidth, height=6.2cm,
    scaled x ticks=false, scaled y ticks=false,
    axis lines*=left, axis line style={draw=linea},
    tick style={draw=none},
    ymajorgrids, grid style={draw=rejilla},
    tick label style={font=\scriptsize\color{tenue}},
    label style={font=\scriptsize\color{tenue}},
    legend style={font=\scriptsize, draw=none, fill=none, legend columns=-1, at={(0,1.02)}, anchor=south west, /tikz/every even column/.append style={column sep=8pt}},
    legend cell align=left,
    every node near coord/.append style={font=\tiny\color{tinta2}},
    /pgf/number format/.cd, use comma, 1000 sep={\,}, min exponent for 1000 sep=4,
  },
}

% Tablas mensuales
\newcolumntype{R}{>{\hspace{6pt}}r}
\newcolumntype{Q}{>{\hspace{3pt}}r}
\newcommand{\hijo}[1]{\hspace{8pt}{\color{tenue}#1}}
\newcommand{\nm}[1]{{\color{rojo}#1}}
\newcommand{\nmh}[1]{{\color{rojosuave}#1}}

\documento{Micros Tempo · Manual del equipo, versión alfa · 24 de noviembre de 2026}

\begin{document}

\begin{titlepage}
\thispagestyle{empty}
\vspace*{1.2cm}
{\small\bfseries\color{marca}\MakeUppercase{Micros Tempo}}\par\vspace{14pt}
{\fontsize{34}{38}\selectfont\bfseries Manual del equipo}\par\vspace{6pt}
{\fontsize{20}{24}\selectfont\color{tinta2} versión alfa}\par\vspace{18pt}
\begin{minipage}{0.84\linewidth}\large\color{tinta2}
Para quienes entran a trabajar entre diciembre de 2026 y abril de 2027: qué hacemos, por qué existe la empresa, qué se espera de cada puesto y cómo se trabaja cada semana.
\end{minipage}
\vfill
\begin{nota}
\textbf{Cómo usar este manual.} Léelo completo en tu primera semana y vuelve a la parte de tu puesto cada vez que tengas una duda. Si algo de lo que dice no coincide con lo que ves en la calle, avísale a Gerson: el manual se corrige con lo que aprendamos juntos.
\end{nota}
\end{titlepage}

{\setcounter{tocdepth}{1}\tableofcontents}
\newpage

% ════════════════════════════════════════════════════════════════════════
\section{Por qué existe Micros Tempo}

\subsection{El problema que resolvemos}

Un chofer de combi no sabe a cuántos minutos va la unidad de adelante ni la de atrás. Ante la duda acelera, porque llegar tarde a un paradero significa que otro se lleva a los pasajeros. Por eso existen los dateros: personas que se paran en una esquina, anotan a qué hora pasó cada unidad y le venden ese dato al chofer o al cobrador que viene detrás, por S/~0,10 a S/~0,50 cada vez.

El dato del datero llega tarde, se da a gritos desde la vereda y solo sirve en la esquina donde está parado. Aun así, los choferes lo pagan todos los días, porque saber dónde va la unidad de adelante cambia la forma de manejar. Ese mercado ya existe, y es la mejor prueba de que lo que vendemos sirve.

\subsection{Por qué importa en Juliaca}

Según la Policía Nacional, entre el 1 de enero y el 3 de mayo de 2026 Juliaca fue la segunda ciudad del país en muertos por accidentes de tránsito, con 96 fallecidos en 74 accidentes. Las autoridades locales atribuyen entre el 80\,\% y el 85\,\% de los accidentes al factor humano, sobre todo al exceso de velocidad.

Un estudio hecho en Santiago de Chile comparó a choferes que cobraban por pasajero con choferes de sueldo fijo en las mismas calles: los primeros tuvieron cerca de un 67\,\% más de accidentes por kilómetro. Ese estudio habla de cómo se les paga a los choferes y no de nuestra app, así que nunca lo usamos para prometer que la app reduce accidentes.

\subsection{Lo que hace la app}

Micros Tempo le muestra al chofer, cada tres segundos y en números grandes, cuántos minutos le lleva la unidad de adelante y cuántos la de atrás. Con ese dato el chofer sabe cuándo no hace falta correr, que es la mayoría de las veces. La empresa de transporte ve toda su flota en un mapa, recibe las emergencias al instante, habla con cada unidad y cuenta las vueltas y las horas de cada persona sin planillas.

\begin{aviso}
\textbf{La frase que resume el producto:} el rastreo mide la vuelta, no a la persona. La app no transmite nada hasta que el chofer desliza «Salir a ruta», el chofer puede marcarse ausente para comer y la app no le reclama nada mientras maneja.
\end{aviso}

% ════════════════════════════════════════════════════════════════════════
\section{Cómo funciona el negocio}

\subsection{Lo que se cobra}

Cada empresa paga S/~0,50 por unidad por día, con el IGV incluido. Una combi paga S/~15 al mes y una ruta de 40 combis paga unos S/~600 al mes. Lo que paga una combi en todo un mes es menos de lo que su chofer le paga al datero en una semana, y por eso el precio rara vez es lo que frena una venta.

\subsection{Por qué tu trabajo es el centro de la empresa}

Atender una ruta más casi no le cuesta nada a la empresa: el servidor de una ruta cuesta menos del 10\,\% de lo que esa ruta paga. Por eso el negocio crece al ritmo en que se firman rutas, y cada ruta que se firma deja dinero todos los meses mientras siga contenta con el servicio. Una ruta que no se firma en un mes no se recupera después: es un mes menos de ingreso para siempre.

\subsection{El plan del primer año}

\begin{center}
\begin{tabularx}{\linewidth}{@{}>{\bfseries}p{3.6cm}X@{}}
\toprule
Diciembre de 2026 & Dos rutas de prueba, que todavía no pagan.\\
Enero a julio de 2027 & El equipo firma hasta 10 rutas por mes y llena Juliaca y Puno, unas 60 rutas, en julio.\\
Agosto a noviembre de 2027 & Pausa: no se firman rutas nuevas y se prepara la entrada a Arequipa o Cusco.\\
Diciembre de 2027 en adelante & Se vuelven a firmar 10 rutas por mes en la ciudad grande.\\
\bottomrule
\end{tabularx}
\end{center}

La empresa gana todos los meses desde mayo de 2027 si cada persona cumple su meta. Si las metas no se cumplen, el negocio sigue siendo viable, pero crece más lento y necesita más capital, así que la meta de cada uno es la parte del plan que más pesa.

% ════════════════════════════════════════════════════════════════════════
\section{Quién es quién}

\subsection{En Micros Tempo}

\begin{center}
\begin{tabularx}{\linewidth}{@{}>{\bfseries}p{3.6cm}p{3.4cm}X@{}}
\toprule
\normalfont Persona & Desde & Su trabajo\\
\midrule
Gerson & Diciembre de 2026 & Es el dueño. Vende hasta febrero y desde marzo se encarga del sistema: el servidor, la app y la seguridad.\\
Vendedor A & Diciembre de 2026 & Medio tiempo en diciembre y tiempo completo desde enero. Vende e instala rutas.\\
Vendedor B & Febrero de 2027 & Medio tiempo. Vende a empresas pequeñas y cuida a las rutas activas.\\
Vendedor D & Marzo de 2027 & Tiempo completo. Recibe la cartera de Gerson y abre Puno.\\
Asistente de oficina & Abril de 2027 & Medio tiempo. Factura, cobra, atiende el teléfono y archiva.\\
Contador externo & Diciembre de 2026 & Lleva los impuestos y la planilla.\\
\bottomrule
\end{tabularx}
\end{center}

\subsection{En cada empresa de transporte}

\begin{itemize}
\item \textbf{La gerencia} decide si contrata el servicio y firma el contrato. En muchas empresas la decisión pasa por la asamblea de socios, porque los socios son dueños de las combis.
\item \textbf{El despacho} es la persona o el equipo que controla la ruta desde una computadora. Da de alta a los choferes y cobradores, cambia claves, atiende las emergencias y es el primer nivel de soporte. Un despacho bien capacitado es la mitad del éxito de una ruta.
\item \textbf{El chofer} lleva la app en un celular Android en soporte y es el único teléfono de la combi que manda la posición.
\item \textbf{El cobrador} usa la app en modo acompañante: ve la misma brecha y el mismo mapa, pero su celular no manda la posición.
\end{itemize}

% ════════════════════════════════════════════════════════════════════════
\section{Tu puesto: objetivos y razones}

Cada puesto tiene una meta que se puede contar y una razón que explica por qué esa meta y no otra. Si un día no sabes qué hacer primero, vuelve a la razón de tu puesto.

\begin{rol}{Vendedor A · tiempo completo}
{\small\color{tenue}Medio tiempo en diciembre de 2026 (S/~700) y tiempo completo desde enero de 2027 (S/~1400) · S/~200 por ruta firmada}\par\vspace{4pt}
\textbf{Objetivo.} Firmar 1 ruta en enero y 4 rutas por mes desde febrero, con la app funcionando en al menos 32 de las 40 unidades de cada ruta.

\textbf{La razón.} Una persona a tiempo completo alcanza a firmar una ruta por semana, porque cada ruta exige conversar con la gerencia, a veces con la asamblea, e instalar la app en decenas de teléfonos. El primer mes se cuenta 1 ruta porque es el mes de aprender el producto y armar la lista de empresas.

\textbf{En diciembre.} Acompaña las 2 rutas de prueba, anota lo que piden los choferes y el despacho, y arma con Gerson la lista de empresas de Juliaca.
\end{rol}

\begin{rol}{Vendedor B · medio tiempo}
{\small\color{tenue}Desde febrero de 2027 · S/~700 · S/~200 por ruta firmada}\par\vspace{4pt}
\textbf{Objetivo.} Firmar 1 ruta en febrero y 2 por mes desde marzo, y recuperar las rutas activas que dejan de usar la app.

\textbf{La razón.} Medio tiempo alcanza para una ruta cada dos semanas. Además, una ruta que deja de usar la app termina dejando de pagar, y recuperarla cuesta menos que firmar una nueva; por eso B visita a las rutas que tienen menos del 80\,\% de sus unidades reportando.
\end{rol}

\begin{rol}{Vendedor D · tiempo completo}
{\small\color{tenue}Desde marzo de 2027 · S/~1400 · S/~200 por ruta firmada}\par\vspace{4pt}
\textbf{Objetivo.} Firmar 1 ruta en marzo y 4 por mes desde abril, abrir Puno y, durante la pausa, dejar listos 10 contratos en la ciudad grande.

\textbf{La razón.} D entra para que Gerson deje de vender y se dedique al sistema, que para entonces atiende cientos de combis. En su primer mes, Gerson lo acompaña a sus reuniones y le pasa sus contactos, para que D no empiece de cero.
\end{rol}

\begin{rol}{Asistente de oficina · medio tiempo}
{\small\color{tenue}Desde abril de 2027 · S/~700 · no vende}\par\vspace{4pt}
\textbf{Objetivo.} Emitir todas las facturas en los tres primeros días del mes y tener cobrado el 95\,\% al día 15 del mes siguiente.

\textbf{La razón.} La empresa paga sueldos e impuestos con lo que cobra, y una factura que se emite tarde se cobra tarde. Además, el IGV se le paga a SUNAT cuando se emite la factura, aunque el cliente todavía no haya pagado, así que cada día de atraso en la cobranza lo financia la empresa.

\textbf{También.} Atiende las llamadas de las empresas, deriva al despacho de cada empresa lo que el despacho puede resolver, archiva los contratos y los consentimientos firmados, y le prepara al contador la información del mes.
\end{rol}

\subsection{Cuándo una ruta cuenta como firmada}

Una ruta cuenta para tu meta y para tu comisión cuando se cumplen tres cosas: la empresa firmó el contrato, el despacho está capacitado y al menos 32 de sus 40 unidades reportan su posición durante una semana. La razón es simple: una ruta con contrato y sin la app funcionando no paga, o deja de pagar al segundo mes.

\subsection{Tu pago}

El sueldo se paga a fin de cada mes. La comisión de S/~200 por ruta se paga con el sueldo del mes en que la ruta cuenta como firmada. Las rutas que trae Gerson no generan comisión, y las 2 rutas de prueba de diciembre sí la generan para quien las acompañó.

% ════════════════════════════════════════════════════════════════════════
\section{Cómo se vende una ruta, paso a paso}

Cada paso tiene un resultado que se puede comprobar. No pases al siguiente sin tener el resultado del anterior.

\begin{enumerate}[label=\textbf{\arabic*.}, leftmargin=18pt, itemsep=5pt]
\item \textbf{Buscar la empresa.} Arma una lista con el nombre de la empresa, su ruta, cuántas unidades tiene, quién es el gerente y cómo ubicarlo. El resultado es una cita con la gerencia.
\item \textbf{Primera visita a la gerencia.} Pregunta antes de mostrar: cuántos dateros pagan sus choferes, cómo controlan hoy las vueltas y qué pasa cuando hay una emergencia en la ruta. El resultado es saber qué le duele a esa empresa y quién decide.
\item \textbf{Demostración.} Muestra la app del chofer y el panel del despacho con los datos de la demostración, nunca con los datos de otra empresa. Usa el manual de usuario de Micros Tempo como guía. El resultado es una fecha para presentar el servicio a los socios, si hace falta.
\item \textbf{Asamblea de socios.} Muchos socios son dueños de su combi y temen que la app sirva para vigilarlos. Explica lo que la app no hace (sección 6) antes de que te lo pregunten. El resultado es el acuerdo de contratar.
\item \textbf{Contrato.} La empresa firma el contrato del servicio y el contrato de encargo de datos personales, que dice que los datos son de la empresa y que Micros Tempo solo los guarda por su cuenta. El resultado son los dos contratos firmados y entregados a la oficina.
\item \textbf{Consentimientos.} Cada chofer y cada cobrador lee y firma el consentimiento de datos personales; la empresa guarda el original y le da una copia. El resultado es un consentimiento firmado por cada persona que va a usar la app.
\item \textbf{Instalación.} Instala la app Android en el celular de cada chofer, dale todos los permisos de ubicación y quita la restricción de batería; en Xiaomi, Huawei y Oppo activa también el inicio automático. Sin esos ajustes, el celular deja de mandar la posición con la pantalla apagada. El resultado es cada unidad reportando en el mapa del despacho.
\item \textbf{Capacitación del despacho.} Enséñale al despacho a dar de alta personas, cambiar claves, atender un SOS y leer los números de la ruta. El resultado es un despacho que resuelve solo los problemas de todos los días.
\item \textbf{Primer mes.} Visita la ruta en la primera y en la tercera semana, revisa qué unidades no reportan y resuélvelo con el despacho. El resultado es una ruta que paga su segunda factura completa y a tiempo.
\end{enumerate}

% ════════════════════════════════════════════════════════════════════════
\section{Lo que nunca se promete}

Un chofer desconfiado comprueba cada promesa en su primer día con la app. Una promesa falsa le quita la confianza a toda la ruta, así que la regla es esta: los números de afuera describen el problema, y la app se describe por lo que hace.

\begin{center}
\begin{tabularx}{\linewidth}{@{}XX@{}}
\toprule
\textbf{Lo que sí se dice} & \textbf{Lo que nunca se dice}\\
\midrule
«Juliaca fue la segunda ciudad del país en muertos por accidentes de tránsito entre enero y mayo de 2026.» & «Esta app reduce los accidentes.» Nadie lo ha medido.\\
«La app te dice cuándo no hace falta correr.» & «La app no distrae.» Cualquier pantalla en un vehículo distrae algo.\\
«El mapa se guarda en el teléfono y no se vuelve a bajar.» & «La app funciona sin señal.» Sin señal, el mapa sigue ahí, pero la brecha se congela hasta que vuelve la señal.\\
«Un turno gasta unos 40 MB de datos con 20 unidades en la ruta.» & Un ahorro exacto de datos o de dinero, porque depende de cada chofer.\\
«Con la app Android, la posición sigue llegando con la pantalla apagada.» & Que funciona igual en cualquier celular: necesita Android, permisos de ubicación y sin restricción de batería.\\
\bottomrule
\end{tabularx}
\end{center}

Hay una limitación que tienes que decir siempre: un SOS o un mensaje no le suena al despacho si el panel no está abierto. Por eso el despacho tiene que trabajar con el panel abierto durante todo el turno.

% ════════════════════════════════════════════════════════════════════════
\section{Lo que te van a preguntar}

\begin{center}
\begin{tabularx}{\linewidth}{@{}>{\itshape}p{4.6cm}X@{}}
\toprule
\normalfont\textbf{Te dicen} & \textbf{Respondes}\\
\midrule
«Me van a estar mirando todo el día.» & La app no manda nada hasta que deslizas «Salir a ruta». Cuando paras a comer, te marcas ausente y sales de la medición, y si te olvidas de volver, la app deja de transmitir sola para no mandar la ubicación de tu casa.\\
«Si me salgo de la ruta me van a sancionar.» & El sistema no te dice nada mientras manejas y solo marca un desvío desde los 300 metros, que son tres cuadras: esquivar un embotellamiento no es salirse de la ruta.\\
«Esto es para que la empresa nos apure más.» & Es al revés: la brecha te dice cuándo no hace falta correr. Hoy aceleras sin saber si estás atrasado.\\
«Me va a gastar los datos.» & Un turno gasta unos 40 MB con 20 unidades en la ruta, y el mapa se guarda en el teléfono. Si eres cobrador, tu celular ni siquiera manda la posición.\\
«¿Y si no tengo señal?» & La app abre con lo último que guardó y el mapa está en el teléfono. Lo que se congela es la brecha, hasta que vuelva la señal.\\
«Otro sistema más que nos van a descontar.» & El precio lo decide la empresa con sus socios. Del lado del chofer, lo que se agrega es su propio registro de horas y vueltas, que hoy no tiene en ninguna parte.\\
\bottomrule
\end{tabularx}
\end{center}

% ════════════════════════════════════════════════════════════════════════
\section{Datos personales y confidencialidad}

Los datos que mueve la app son de cada empresa de transporte y de cada chofer: dónde estuvo, cuántas horas trabajó y qué escribió en el chat. La Ley 29733 de Protección de Datos Personales obliga a cuidarlos, y la confianza de las empresas depende de que nadie del equipo los use para otra cosa.

\begin{itemize}
\item Nunca mires la ubicación, los mensajes o las horas de una persona por curiosidad o para un asunto personal. Solo se miran para resolver un problema que la empresa te pidió resolver.
\item Nunca muestres los datos de una empresa a otra empresa. En las demostraciones se usa solo la empresa de prueba.
\item Nunca compartas tu clave ni pidas la clave de un chofer. Si alguien olvida su clave, la cambia el despacho de su empresa.
\item Si ves algo raro, como una persona que entra con la cuenta de otra o datos que no deberías ver, avísale a Gerson ese mismo día. El contrato nos da 24 horas para avisarle a la empresa, y la ley le da 48 horas a la empresa para avisar a la autoridad, contadas desde que se supo del problema.
\end{itemize}

% ════════════════════════════════════════════════════════════════════════
\section{Cómo medimos el trabajo}

Cada viernes, cada vendedor actualiza el tablero del equipo con cuatro números de su semana. Los números no sirven para castigar a nadie: sirven para ver a tiempo dónde se traba una venta y ayudar antes de que se pierda el mes.

\begin{center}
\begin{tabularx}{\linewidth}{@{}>{\bfseries}p{4.4cm}X@{}}
\toprule
Reuniones con gerencias & Cuántas empresas nuevas escucharon la presentación.\\
Contratos firmados & Cuántas empresas firmaron los dos contratos.\\
Unidades reportando & Cuántas unidades de tus rutas nuevas mandan su posición.\\
Rutas en riesgo & Cuántas rutas tuyas tienen menos del 80\,\% de sus unidades reportando.\\
\bottomrule
\end{tabularx}
\end{center}

Una vez al mes, Gerson revisa con cada uno sus números frente a la meta. Si un vendedor no llega a su meta dos meses seguidos, se busca la causa con él antes de cambiar nada: el problema puede ser el argumento, la lista de empresas o un paso del proceso que no funciona.

% ════════════════════════════════════════════════════════════════════════
\section{La pausa de agosto a noviembre de 2027}

Cuando Juliaca y Puno lleguen a unas 60 rutas, ya no quedará a quién venderle en esas dos ciudades. Desde agosto, el equipo deja de firmar rutas y se dedica cuatro meses a preparar la entrada a Arequipa o Cusco, que tienen muchas más rutas.

\begin{itemize}
\item \textbf{En agosto,} los vendedores le proponen a Gerson la ciudad grande, con la lista de empresas y sus rutas.
\item \textbf{En septiembre y octubre,} el equipo consigue 30 reuniones con gerencias de la ciudad elegida.
\item \textbf{En noviembre,} deja 10 contratos listos para firmar en diciembre.
\item \textbf{Durante toda la pausa,} B y la asistente de oficina cuidan que las rutas de Juliaca y Puno sigan reportando y pagando a tiempo.
\end{itemize}

La pausa no es tiempo libre: es el trabajo que decide cuánto crece la empresa en 2028.

% ════════════════════════════════════════════════════════════════════════
\section{Tu primera semana}

\begin{enumerate}[label=\textbf{Día \arabic*.}, leftmargin=40pt]
\item Lee este manual y el manual de usuario de Micros Tempo, y firma tu contrato de trabajo.
\item Instala la app en tu celular, entra con una cuenta de prueba y recorre todas las pantallas del chofer.
\item Pasa una mañana con el despacho de una ruta activa y mira cómo usa el panel.
\item Acompaña a Gerson o a un vendedor a una reunión con una gerencia y a una instalación.
\item Arma tu lista de 20 empresas para el mes, con su ruta, sus unidades y su gerente.
\end{enumerate}

Al terminar la semana tienes que poder explicar, sin leer, qué es la brecha, qué hace la app y qué no hace. Si no puedes, repite la presentación con Gerson hasta que salga natural.

% ════════════════════════════════════════════════════════════════════════
\section{Palabras que usamos}

\begin{center}
\begin{tabularx}{\linewidth}{@{}>{\bfseries}p{3.4cm}X@{}}
\toprule
Ruta & El recorrido de una empresa de transporte, con todas las unidades que lo hacen. Cobramos por ruta.\\
Unidad & Cada combi o micro de la ruta. Una ruta típica tiene 40.\\
Brecha & Los minutos que separan a una unidad de la de adelante o de la de atrás. Es el número grande de la app.\\
Pelotón & Varias unidades pegadas, que se reparten a los pasajeros y dejan paraderos vacíos detrás.\\
Datero & La persona que anota en una esquina a qué hora pasa cada unidad y le vende el dato al chofer.\\
Despacho & Quien controla la ruta desde el panel de la empresa y da el primer nivel de soporte.\\
Salir a ruta & El gesto de deslizar con el que el chofer empieza a mandar su posición.\\
Ausente & El estado que marca el chofer cuando para a comer o a reparar algo; sale de la medición.\\
SOS & La alerta de emergencia que el chofer o el cobrador mandan deslizando, sin desbloquear el teléfono.\\
Consentimiento & El papel que firma cada chofer y cobrador para que la empresa trate sus datos.\\
\bottomrule
\end{tabularx}
\end{center}

\end{document}
